\label{sec:methodology}

This study uses an application-oriented experimental methodology. The goal is
to determine which open LLMs are practical for railway-domain use when the
hardware budget is fixed to a single 24GB GPU. The method combines baseline
evaluation, low-bit inference comparison, and QLoRA-based domain adaptation.

\textbf{Baseline evaluation.} Each model is first evaluated without fine-tuning.
This establishes the zero-adaptation capability of the model on public
benchmarks and on the railway-domain QA test set. The local railway evaluation
uses exact matching after answer normalization, while public benchmarks use the
metrics defined by their respective evaluation harness tasks.

\textbf{Efficiency measurement.} Accuracy alone is insufficient under a strict
single-GPU constraint. For each model and precision setting, the pipeline records
peak GPU memory, latency, and generated-token throughput on a fixed efficiency
prompt set. These values are analyzed together with task scores to identify
models that provide a favorable performance-efficiency tradeoff.

\textbf{Railway-domain adaptation.} For selected models, QLoRA is used to adapt
the model to railway-domain samples. The adapted model is represented as a
small adapter rather than a full copy of the base model. After adaptation, the
same railway-domain evaluation is repeated. The adaptation gain is defined as:
\[
\Delta S_{\mathrm{domain}} = S_{\mathrm{after}} - S_{\mathrm{before}},
\]
where \(S_{\mathrm{before}}\) and \(S_{\mathrm{after}}\) are the domain test
scores before and after QLoRA.

\textbf{Deployment-oriented analysis.} The final analysis compares models along
four dimensions: task score, memory footprint, latency, and adaptation response.
This avoids treating the paper as a simple leaderboard. A model with the
highest benchmark score may not be the best practical choice if it is too slow
or too expensive to adapt. Conversely, a smaller model may become attractive if
QLoRA yields strong domain improvement at low cost.
